% ============================================================================
%  附录 C：术语中英对照
% ============================================================================
\chapter{术语中英对照}
\label{app:glossary}

阅读英文文献与面试时，术语的英文表达往往比中文更常用。下表按主题分组，
每条给出英文、中文与一句话解释（含所在章节）。

\section{市场与产品}

{\tablefont
\begin{longtable}{@{}L{4.6cm}L{3.6cm}L{6.6cm}@{}}
\toprule
\textbf{英文} & \textbf{中文} & \textbf{说明与章节} \\
\midrule
\endfirsthead
\toprule
\textbf{英文} & \textbf{中文} & \textbf{说明与章节} \\
\midrule
\endhead
\bottomrule
\endlastfoot
Primary / Secondary Market & 一级 / 二级市场 & 证券首次发行与后续流通市场（第 1 章） \\
Money Market & 货币市场 & 期限不超过一年的短期资金市场（第 1 章） \\
Repurchase Agreement (Repo) & 回购协议 & 以证券为抵押的短期融资，是货币市场核心工具（第 1 章） \\
Commercial Paper & 商业票据 & 企业发行的短期无担保票据（第 1 章） \\
Yield Curve & 收益率曲线 & 不同期限的无风险收益率关系（第 1 章） \\
Duration / Convexity & 久期 / 凸性 & 债券价格对利率的一阶 / 二阶敏感度（第 1、5 章） \\
Immunization & 免疫策略 & 用久期与凸性匹配资产负债的利率风险（第 1 章） \\
DuPont Analysis & 杜邦分解 & 把 ROE 拆成净利率、周转率与权益乘数（第 1 章） \\
Forward / Futures & 远期 / 期货 & 场外与交易所的远期合约（第 4 章） \\
Swap / Swaption & 互换 / 互换期权 & 交换现金流的合约与其期权（第 4 章） \\
Cap / Floor & 利率上限 / 下限 & 一系列利率看涨 / 看跌期权（第 4 章） \\
Autocallable & 自动赎回结构 & 触发障碍即提前赎回的结构化产品（第 4 章） \\
Credit Default Swap (CDS) & 信用违约互换 & 转移参考实体违约风险的合约（第 4 章） \\
Collateralized Debt Obligation & 担保债务凭证 & 把信用风险池分层的结构化产品（第 4 章） \\
Municipal Bond / TEY & 市政债 / 税收等价收益率 & 免税债券与应税债券的收益比较（第 12 章） \\
\end{longtable}
}

\section{定价与随机分析}

{\tablefont
\begin{longtable}{@{}L{4.6cm}L{3.6cm}L{6.6cm}@{}}
\toprule
\textbf{英文} & \textbf{中文} & \textbf{说明与章节} \\
\midrule
\endfirsthead
\toprule
\textbf{英文} & \textbf{中文} & \textbf{说明与章节} \\
\midrule
\endhead
\bottomrule
\endlastfoot
No-Arbitrage & 无套利 & 定价的第一性原理（第 4、6 章） \\
Law of One Price & 一价定律 & 相同现金流必须有相同价格（第 4 章） \\
Equivalent Martingale Measure & 等价鞅测度 & 风险中性测度（第 6 章） \\
Radon--Nikodym Derivative & Radon--Nikodym 导数 & 测度变换的密度过程（第 6 章） \\
Girsanov Theorem & Girsanov 定理 & 漂移变换：改变测度只改漂移不改波动率（第 6、10 章） \\
Num\'eraire & 计价单位 & 用于折现的基准资产（第 6 章） \\
Market Completeness & 市场完备性 & 所有未定权益可复制的条件（第 6、10 章） \\
Feynman--Kac Formula & Feynman--Kac 公式 & 连接 PDE 与期望的桥梁（第 10 章） \\
It\^o's Lemma & 伊藤引理 & 随机微积分的链式法则（第 10 章） \\
Quadratic Variation & 二次变差 & 布朗运动的「时间刻度」（第 10 章） \\
Geometric Brownian Motion & 几何布朗运动 & BSM 的标的假设（第 4、10 章） \\
Ornstein--Uhlenbeck Process & OU 过程 & 均值回复过程，用于利率与配对交易（第 10 章） \\
Jump Diffusion & 跳跃扩散 & 加入跳跃的连续时间模型（第 10 章） \\
L\'evy Process & L\'evy 过程 & 平稳独立增量的通用跳跃过程（第 10 章） \\
Implied / Local Volatility & 隐含 / 局部波动率 & 从期权价格反解 / Dupire 公式推出的波动率（第 4 章） \\
Volatility Smile / Skew & 波动率微笑 / 偏斜 & 隐含波动率随行权价变化的形态（第 4 章） \\
Stochastic Volatility & 随机波动率 & Heston 等模型（第 4、6 章） \\
Variance Swap & 方差互换 & 交易实现方差的衍生品（第 6 章） \\
Fast Fourier Transform Pricing & 傅里叶定价 & 用特征函数一次算出整条期权价格曲线（第 6 章） \\
Least-Squares Monte Carlo & 最小二乘蒙特卡洛 & 美式 / 百慕大期权定价（第 6 章） \\
XVA (CVA/DVA/FVA) & 估值调整 & 交易对手信用、自身信用与资金成本调整（第 4 章） \\
\end{longtable}
}

\section{统计、计量与机器学习}

{\tablefont
\begin{longtable}{@{}L{4.6cm}L{3.6cm}L{6.6cm}@{}}
\toprule
\textbf{英文} & \textbf{中文} & \textbf{说明与章节} \\
\midrule
\endfirsthead
\toprule
\textbf{英文} & \textbf{中文} & \textbf{说明与章节} \\
\midrule
\endhead
\bottomrule
\endlastfoot
Stylized Facts & 程式化事实 & 收益的肥尾、波动率聚集等经验规律（第 9 章） \\
Volatility Clustering & 波动率聚集 & 大波动后跟大波动（第 9 章） \\
Leverage Effect & 杠杆效应 & 负收益对波动率影响更强（第 9 章） \\
Heteroskedasticity & 异方差 & 方差不恒定（第 8、9 章） \\
Autocorrelation (ACF/PACF) & 自相关 / 偏自相关 & 时间序列识别工具（第 9 章） \\
Cointegration & 协整 & 非平稳序列间的长期均衡关系（第 9 章） \\
Error Correction Model & 误差修正模型 & 协整系统的短期动态（第 9 章） \\
State-Space Model & 状态空间模型 & 潜变量与观测方程的框架（第 9 章） \\
Kalman Filter & 卡尔曼滤波 & 状态空间模型的递推估计（第 9 章） \\
Realized Volatility & 已实现波动率 & 高频收益平方和（第 9 章） \\
HAR Model & HAR 模型 & 已实现波动的多尺度（日/周/月）回归（第 9 章） \\
Extreme Value Theory & 极值理论 & 尾部建模的理论基础（第 8 章） \\
Peaks Over Threshold & 超阈值法 & 用 GPD 拟合尾部（第 8 章） \\
Maximum Likelihood / GMM & 极大似然 / 广义矩 & 两类主流估计方法（第 8 章） \\
Shrinkage Estimator & 收缩估计 & 用结构化目标矩阵稳定样本协方差（第 11 章） \\
Bayesian Shrinkage & 贝叶斯收缩 & Black--Litterman 的本质（第 8、11 章） \\
Fama--MacBeth Regression & Fama--MacBeth 两步回归 & 资产定价的标准检验方法（第 8 章） \\
Look-Ahead Bias & 前视偏差 & 使用了未来信息（第 7 章） \\
Survivorship Bias & 幸存者偏差 & 样本只剩存活个体（第 7 章） \\
Data Snooping & 数据窥探 & 反复试探导致假显著（第 7 章） \\
Purged K-Fold & 净化交叉验证 & 带间隔的时间序列交叉验证（第 7 章） \\
Deflated Sharpe Ratio & 通缩夏普比率 & 对多重检验修正后的夏普（第 7 章） \\
Implementation Shortfall & 实施缺口 & 决策价格与最终成交价的差异（第 7 章） \\
\end{longtable}
}

\section{组合、风险与交易}

{\tablefont
\begin{longtable}{@{}L{4.6cm}L{3.6cm}L{6.6cm}@{}}
\toprule
\textbf{英文} & \textbf{中文} & \textbf{说明与章节} \\
\midrule
\endfirsthead
\toprule
\textbf{英文} & \textbf{中文} & \textbf{说明与章节} \\
\midrule
\endhead
\bottomrule
\endlastfoot
Mean--Variance Optimization & 均值--方差优化 & Markowitz 框架（第 11 章） \\
Efficient Frontier & 有效前沿 & 给定风险下收益最高的组合集合（第 11 章） \\
Tangency Portfolio & 切点组合 & 最优风险资产组合（第 11 章） \\
Two-Fund Separation & 两基金分离 & 所有有效组合可由两个基金张成（第 11 章） \\
Risk Parity / ERC & 风险平价 / 等风险贡献 & 让各资产风险贡献相等（第 11 章） \\
Risk Budgeting & 风险预算 & 指定各资产的风险贡献比例（第 11 章） \\
Hierarchical Risk Parity & 层次风险平价 & 用聚类替代矩阵求逆（第 11 章） \\
Maximum Diversification & 最大分散度 & 最大化分散度比率（第 11 章） \\
Kelly Criterion & Kelly 准则 & 最大化对数增长率的仓位（第 11 章） \\
Value at Risk / Expected Shortfall & 在险价值 / 预期亏空 & 分位数与尾部均值风险度量（第 4 章） \\
Coherent Risk Measure & 一致性风险度量 & 满足次可加性等公理的度量（第 4 章） \\
Backtesting & 回测 / 风险回测 & 检验策略或风险模型的样本外表现（第 7、9 章） \\
Kupiec / Christoffersen Test & 覆盖率 / 独立性检验 & VaR 回测的两种标准检验（第 9 章） \\
Performance Attribution & 绩效归因 & 把超额收益分解为配置与选股（第 11 章） \\
Information Coefficient / Ratio & 信息系数 / 信息比率 & 信号质量与风险调整后主动收益（第 7 章） \\
Fundamental Law of Active Management & 主动管理基本定律 & $IR\approx IC\sqrt{\text{Breadth}}$（第 7 章） \\
Limit Order Book & 限价订单簿 & 微观结构的核心数据结构（第 7 章） \\
Order Flow Imbalance & 订单流不平衡 & 买卖压力的度量（第 7 章） \\
Market Impact & 市场冲击 & 交易对价格的反作用（第 7、11 章） \\
Adverse Selection & 逆向选择 & 做市商面对知情交易者的损失（第 7 章） \\
\end{longtable}
}

\section{金融科技、Web3 与监管}

{\tablefont
\begin{longtable}{@{}L{4.6cm}L{3.6cm}L{6.6cm}@{}}
\toprule
\textbf{英文} & \textbf{中文} & \textbf{说明与章节} \\
\midrule
\endfirsthead
\toprule
\textbf{英文} & \textbf{中文} & \textbf{说明与章节} \\
\midrule
\endhead
\bottomrule
\endlastfoot
Open Banking & 开放银行 & 通过 API 开放银行数据（第 2 章） \\
Weight of Evidence / IV & WOE / 信息量 & 评分卡特征变换与筛选（第 2 章） \\
Expected Loss & 预期损失 & $EL=PD\times LGD\times EAD$（第 2 章） \\
RAROC & 风险调整资本收益率 & 用经济资本衡量的收益率（第 2 章） \\
Model Risk Management & 模型风险治理 & 美联储 SR 11-7 框架（第 2 章） \\
RegTech / SupTech & 监管科技 & 用技术降低合规成本（第 2 章） \\
Automated Market Maker & 自动做市商 & 用不变量定价的链上做市（第 3 章） \\
Constant Product Formula & 恒定乘积公式 & $xy=k$（第 3 章） \\
Impermanent Loss & 无常损失 & 做市相对持有的机会成本（第 3 章） \\
Total Value Locked & 总锁仓价值 & 协议内锁定的资产规模（第 3 章） \\
Maximal Extractable Value & 最大可提取价值 & 排序交易获取的额外收益（第 3 章） \\
Rollup (Optimistic / ZK) & 二层汇总 & 把执行搬到链下的扩容方案（第 3 章） \\
Reentrancy Attack & 重入攻击 & 智能合约经典漏洞（第 3 章） \\
Tokenomics & 代币经济学 & 代币供给、激励与价值捕获设计（第 3 章） \\
Howey Test & Howey 测试 & 判断是否构成证券的美国判例标准（第 3 章） \\
MiCA & 加密资产市场监管法案 & 欧盟加密资产监管框架（第 3 章） \\
Know Your Customer & 了解你的客户 & 客户身份识别义务（第 2、12 章） \\
Anti-Money Laundering & 反洗钱 & 可疑交易监测与报告（第 12 章） \\
Fiduciary Duty & 受托责任 & 投资顾问的最高行为标准（第 12 章） \\
Suitability / Reg BI & 适合性 / 最佳利益规则 & 向客户推荐产品时的义务（第 12 章） \\
\end{longtable}
}
